% from ega3-1-fr.tex
\oldpage[III]{129}

quer (0, 13.3.1), et l'homomorphisme canonique
$H^n(\widehat X,\widehat{\mathcal F})\to
\varprojlim_kH^n(X,\mathcal F_k)$ est bijectif pour tout $n\geq0$.

Comme le système projectif $(\mathrm H/\mathrm R_k)_{k\geq0}$ est strict,
on peut passer à la limite projective dans les suites exactes
\[
0\longrightarrow\mathrm H/\mathrm R_k\longrightarrow\mathrm H_k
\longrightarrow\mathrm Q_k\longrightarrow0
\tag{4.1.7.22}\label{III.4.1.7.22-fr}
\]
(0, 13.2.2)~; comme $\nu_m(\mathrm Q_{k+m})=0$, on a
$\varprojlim_k\mathrm Q_k=0$, d'où un isomorphisme topologique
\[
\varprojlim_k(\mathrm H/\mathrm R_k)\xrightarrow{\sim}
\varprojlim_k\mathrm H_k.
\]
Mais comme la filtration $(\mathrm R_k)$ de $\mathrm H$ est
$\mathfrak J$-bonne, elle définit sur $\mathrm H$ la topologie
$\mathfrak J$-préadique~; donc $\varprojlim_k(\mathrm H/\mathrm R_k)$ est
le séparé complété de $\mathrm H$ pour la topologie
$\mathfrak J$-préadique, ce qui achève de démontrer
(\hyperref[III.4.1.7-fr]{4.1.7}).
\end{env}
\end{proof}

\begin{env}[4.1.8]
\label{III.4.1.8-fr}
Passons enfin à la démonstration de (\hyperref[III.4.1.5-fr]{4.1.5})~:
pour tout ouvert affine $V$ de $Y$,
$\Gamma(V,\widehat{R^nf_*(\mathcal F)})$ est le séparé complété de
$\Gamma(V,R^nf_*(\mathcal F))$ pour la topologie $\mathfrak J$-préadique
(si $\mathcal J|V=\widetilde{\mathfrak J}$) puisque
$R^nf_*(\mathcal F)$ est un $\mathcal O_Y$-Module cohérent (I, 10.8.4),
et
\[
\Gamma\bigl(V,\varprojlim_kR^nf_*(\mathcal F_k)\bigr)
=\varprojlim_k\Gamma(V,R^nf_*(\mathcal F_k))
\]
(0\textsubscript{I}, 3.2.6)~; le fait que $\varphi_n$ soit un isomorphisme
topologique résulte alors de (\hyperref[III.4.1.7-fr]{4.1.7}) et de
(\hyperref[III.1.4.11-fr]{1.4.11}). D'autre part (toujours en vertu de
(\hyperref[III.1.4.11-fr]{1.4.11})), il résulte de
(\hyperref[III.4.1.7-fr]{4.1.7}) que l'homomorphisme $\psi_{n,V}$ de
(\hyperref[III.4.1.3.3-fr]{4.1.3.3}) est un isomorphisme, donc $\psi_n$
est un isomorphisme par définition de
$R^n\widehat f_*(\widehat{\mathcal F})$.
\end{env}

\begin{corollary}[4.1.9]
\label{III.4.1.9-fr}
Sous les hypothèses de (\hyperref[III.4.1.4-fr]{4.1.4}), pour tout ouvert
affine $V$ de $Y$, l'homomorphisme canonique
\[
H^n(\widehat X\cap f^{-1}(V),\widehat{\mathcal F})
\longrightarrow\Gamma(\widehat Y\cap V,
R^n\widehat f_*(\widehat{\mathcal F}))
\]
est bijectif.
\end{corollary}

\begin{env}[4.1.10]
\label{III.4.1.10-fr}
\emph{Remarque.}
Soit $f:X\to Y$ un morphisme de type fini de préschémas (usuels)
noethériens, et soit $\mathcal F$ un $\mathcal O_X$-Module cohérent dont le
support soit propre sur $Y$ (II, 5.4.10). On sait alors
(\hyperref[III.3.2.4-fr]{3.2.4}) que $R^nf_*(\mathcal F)$ est un
$\mathcal O_Y$-Module cohérent pour tout $n\geq0$. En outre, on peut
toujours supposer que $\mathcal F=u_*(\mathcal G)$, où
$\mathcal G=u^*(\mathcal F)$ est un $\mathcal O_Z$-Module cohérent, $Z$
désignant un sous-préschéma fermé convenable de $X$ dont l'espace
sous-jacent est $\operatorname{Supp}(\mathcal F)$, et $u:Z\to X$
l'injection canonique (I, 9.3.5). Si on pose
\[
\mathcal G_k=\mathcal G\otimes_{\mathcal O_Y}
(\mathcal O_Y/\mathcal J^{k+1}),
\]
on a $\mathcal G_k=u^*(\mathcal F_k)$,
$R^nf_*(\mathcal F_k)=R^n(f\circ u)_*(\mathcal G_k)$ et
$R^nf_*(\mathcal F)=R^n(f\circ u)_*(\mathcal G)$
(\hyperref[III.1.3.4-fr]{1.3.4}), et enfin, compte tenu de (I, 10.9.5),
\[
R^n\widehat f_*(\widehat{\mathcal F})
=R^n\widehat{(f\circ u)}_*(\widehat{\mathcal G}).
\]
On peut alors appliquer (\hyperref[III.4.1.5-fr]{4.1.5}) à $\mathcal G$ et
au morphisme propre $f\circ u$, et on en conclut que sous ces hypothèses,
les résultats de (\hyperref[III.4.1.5-fr]{4.1.5}) sont valables pour
$\mathcal F$ et $f$.
\end{env}

\subsection{Cas particuliers et variantes}
\label{subsection:III.4.2-fr}

La forme la plus utile du th. de comparaison
(\hyperref[III.4.1.5-fr]{4.1.5}) est la suivante~:

\begin{proposition}[4.2.1]
\label{III.4.2.1-fr}
Soient $Y$ un préschéma localement noethérien, $f:X\to Y$ un morphisme
propre, $\mathcal F$ un $\mathcal O_X$-Module cohérent. Alors, pour tout
$y\in Y$ et tout $p$, $(R^pf_*(\mathcal F))_y$ est

